% !TEX root =  main.tex
\chapter{Conclusion}
\label{chap:conclusion}
\pagestyle{plain}

\section{Summary of Contributions}

This thesis presented a comprehensive framework for creating and benchmarking learned database components using large language models and active learning. We demonstrated that large language models can generate diverse, realistic SQL workloads at scale, reducing reliance on manual query collection and synthetic templates. By incorporating uncertainty sampling and Monte Carlo dropout strategies into the training process, we achieved notable improvements in sample efficiency compared to random labeling approaches, enabling the model to reach competitive accuracy levels with a substantially smaller labeling budget. The implemented system provides a complete, schema-agnostic solution that integrates query generation, multi-layer validation, database labeling, feature extraction, and model training into a unified and fully automated pipeline. Through comprehensive partitioning strategies including complexity-based, schema-aware, and uncertainty-based methods, the framework enables efficient learning across different workload characteristics. The extensive experimental evaluation on the IMDB dataset established rigorous evaluation protocols using Q-error metrics and learning curves, providing a reliable basis for comparing different training strategies.


\section{Impact and Implications}

The work presented in this thesis advances the field of learned database systems by addressing the critical challenge of training data generation. Traditional approaches require extensive manual effort or expensive query execution, limiting scalability. Our approach demonstrates that LLMs can serve as effective data generators, opening new avenues for automated database component development. From a practical standpoint, the framework enables more efficient training of learned cardinality estimators for better query optimization, facilitates comprehensive testing of database systems through automated generation of diverse workloads, and provides an open platform for researchers to experiment with learned database components. The reduced training costs also make learned approaches more viable for adoption in commercial database systems.

Beyond database systems, the integration of active learning with generative AI represents a broader methodological contribution. It demonstrates the potential of LLMs in generating structured, constraint-satisfying data, shows how uncertainty-based sampling can be combined effectively with modern neural architectures, and provides a template for building automated machine learning pipelines with minimal human intervention.


\section{Limitations and Future Work}

While the framework achieves significant improvements, several limitations remain. The current implementation focuses on relatively simple conjunctive queries; extending to complex analytical queries with subqueries, aggregations, and window functions requires additional research. Although the pipeline is schema-agnostic by design, performance may vary across different database schemas and data distributions. The Monte Carlo dropout inference and bitmap generation introduce computational overhead that may limit real-time applications. Additionally, timeout handling and error recovery during labeling may introduce noise in the training data through fallback cardinality assignments.

Several promising directions for future work emerge from this thesis. On the query generation side, integrating LLMs with other generative models could enable more complex query patterns, and enhancing LLMs with deeper database semantics could improve the quality of generated workloads. Cross-domain transfer methods could enable learned components trained on one schema to be applied to others.

For active learning, multi-objective optimization that balances uncertainty, diversity, and computational cost in acquisition functions represents a natural extension. Adapting the framework for online active learning in streaming query workloads would enable deployment in production systems, and Bayesian optimization methods could provide more sophisticated uncertainty estimation.

From a systems perspective, integrating learned components into existing database systems with minimal overhead, scaling the pipeline for distributed training over large databases, and developing self-tuning systems that adapt to changing workloads automatically are all important directions. On the theoretical side, providing convergence guarantees for active learning in cardinality estimation, establishing sample complexity bounds, and developing generalization theory for learned estimators would strengthen the foundations of this approach.


\section{Final Remarks}

The research presented in this thesis demonstrates the feasibility and benefits of using large language models for automated training data generation in learned database systems. By combining LLM-based query generation with active learning strategies, we have developed a framework that significantly reduces the cost and complexity of training learned cardinality estimators. The empirical results on the IMDB dataset show that our approach achieves competitive accuracy with substantially fewer labeled examples compared to traditional supervised methods.

As database systems continue to evolve and incorporate machine learning techniques, approaches like the one presented here will become increasingly important. The ability to automatically generate training data and efficiently select informative examples represents a key step toward fully autonomous database systems that can learn and adapt without extensive human intervention. The integration of generative AI with traditional database techniques opens exciting possibilities for the future of data management, and we hope this work inspires further research in this promising direction.